\chapter{From Fourier to the Quantum Fourier Transform}\label{ch:fourier}

\begin{etymology}
\textbf{Period} comes from the Greek \emph{periodos}: \emph{peri-}, ``around'',
and \emph{hodos}, ``way''. A way around, a circuit that comes back to where it
started.

\textbf{Frequency} comes from the Latin \emph{frequens}, ``crowded'',
``repeated''. How often something comes back.
\end{etymology}

Shor's algorithm (Chapter~\ref{ch:encoding}) ends with a quantum Fourier transform
that turns a hidden period into a measurable number. To see why that works, you
have to see what the ordinary Fourier transform does. It was invented to study heat,
and it turned out to be the natural tool for anything periodic.

\section{The problem that started it}

\begin{history}{1807--1822 \textperiodcentered{} Grenoble and Paris}
Joseph Fourier, prefect of the Isère department under Napoleon, studied how heat
spreads through a solid. The equation he wrote had simple solutions shaped like
sines and cosines, but real temperature profiles are not sine-shaped. He claimed
that any function can be written as a sum of sines and cosines, and submitted the
memoir to the Institut de France in 1807. Lagrange objected: a sum of smooth
oscillations, he argued, could not produce a function with corners. The claim was
not fully rigorous, and Peter Gustav Lejeune Dirichlet only proved in 1829 under
which conditions it holds. Fourier published his \emph{Théorie analytique de la
chaleur} in 1822. The idea survived because it worked.
\end{history}

Decomposing a signal into frequencies separates its scales. Each frequency can be
handled on its own and the results put back together, which is often much simpler
than handling the whole signal at once. The central question becomes: how much of
each frequency is in this signal?

\section{Fourier series}

A periodic function with period $T$ can be written as
\[
  f(t) = \sum_{k=-\infty}^{\infty} c_k\, e^{i 2\pi k t / T},
  \qquad
  c_k = \frac{1}{T}\int_0^T f(t)\, e^{-i 2\pi k t/T}\, dt .
\]
Each $c_k$ is how much of the frequency $k/T$ is present in $f$. The exponentials
are used instead of separate sines and cosines because each one packs both into a
single rotation (Chapter~\ref{ch:hilbert}).

This is linear algebra in disguise. Functions form a vector space with the inner
product $\braket{f|g} = \int_0^T \overline{f(t)}\, g(t)\, dt$, and in that space the
exponentials $e^{i2\pi kt/T}$ are orthonormal: the inner product of two different
ones is zero. They are a basis, and the $c_k$ are the coordinates of $f$ in it.
Computing $c_k$ is projecting $f$ onto one pure oscillation.

\section{The discrete Fourier transform}

A computer has no continuous functions, only $N$ samples $x_0, \dots, x_{N-1}$. The
\emph{discrete Fourier transform} (DFT) maps them to $N$ complex coefficients:
\[
  X_k = \sum_{j=0}^{N-1} x_j\, \omega_N^{-jk}, \qquad \omega_N = e^{i 2\pi / N},
\]
and the inverse brings them back,
$x_j = \frac1N\sum_k X_k\,\omega_N^{jk}$. $X_k$ measures how much the sequence
correlates with the oscillation of frequency $k/N$.

For $N = 4$ and the sequence $[1, 0, 1, 0]$, the transform is $[2, 0, 2, 0]$. The
sequence repeats every 2 samples, and the spectrum has weight only at $k = 0$ and
$k = 2$.

\section{The DFT is a unitary matrix}

The DFT is linear, so it is a matrix. With the normalization chosen to make it
symmetric,
\[
  (F_N)_{kj} = \frac{1}{\sqrt N}\, \omega_N^{jk},
\]
and $F_N$ is unitary: $F_N^\dagger F_N = I$. It is a change of basis that preserves
length, from the standard basis to the \emph{Fourier basis}, whose vectors are pure
frequencies.

For $N = 2$,
\[
  F_2 = \frac{1}{\sqrt2}\begin{pmatrix} 1 & 1 \\ 1 & -1 \end{pmatrix} = H .
\]
The Hadamard gate is the Fourier transform of one qubit. The first gate in almost
every quantum circuit is already a Fourier transform. That is why $H\ket0$ is an
even superposition: the state $\ket0$ is a single spike, and the transform of a
spike is flat.

\begin{history}{1965 \textperiodcentered{} IBM and Princeton}
Multiplying by $F_N$ directly costs $O(N^2)$ operations. In 1965 James Cooley and
John Tukey published an algorithm that exploits the recursive structure of $F_N$ to
do it in $O(N\log N)$: the fast Fourier transform. Richard Garwin had pushed for it
because detecting Soviet nuclear tests from seismometer data needed fast spectral
analysis. Gauss had found essentially the same method around 1805, in an unpublished
notebook on asteroid orbits.
\end{history}

\section{The DFT finds periods}

This is the property that matters for the rest of the book. If a sequence of length
$N$ has period $r$, and $r$ divides $N$, its transform is zero everywhere except at
multiples of $N/r$. The $N/r$ copies of the period add up in phase at those
frequencies and cancel exactly at all the others. It is interference, the same
mechanism as in Chapter~\ref{ch:qubit}.

Take $N = 8$ and a spike every 4 samples:
\[
  [1, 0, 0, 0, 1, 0, 0, 0] \;\xrightarrow{\ \text{DFT}\ }\; [2, 0, 2, 0, 2, 0, 2, 0].
\]
The peaks are at $0, 2, 4, 6$, the multiples of $8/4 = 2$. The
spectrum of a periodic signal is sparse, and the spacing of its peaks reveals the
period.

Apply this to $f(a) = 2^a \bmod 15$, which has period 4 (Chapter~\ref{ch:factoring}), and the
transform of its first eight values has peaks at $k = 2, 4, 6$ (the constant part
sits at $k = 0$). Spacing 2, so $r = 8/2 = 4$.

\section{The quantum Fourier transform}

The quantum Fourier transform (QFT) is the same matrix $F_N$, with $N = 2^n$, built as
a circuit on $n$ qubits:
\[
  \mathrm{QFT}\,\ket{j} = \frac{1}{\sqrt N}\sum_{k=0}^{N-1} \omega_N^{jk}\,\ket{k}.
\]
Applied to a superposition $\sum_j x_j\ket j$, it produces $\sum_k X_k\ket k$: the DFT
of the amplitudes.

\begin{history}{1994 \textperiodcentered{} IBM}
Shor's first construction needed a Fourier transform over a large range of values,
and Don Coppersmith, at IBM, showed in 1994 how to build it from Hadamards and
small controlled phase rotations. Coppersmith's report was titled ``An approximate
Fourier transform useful in quantum factoring''. He also showed that the tiniest
rotations can be dropped with little loss of accuracy.
\end{history}

The circuit uses $O(n^2)$ gates: one Hadamard per qubit, controlled rotations
$R_k = \mathrm{diag}(1, e^{i2\pi/2^k})$ between pairs, and swaps at the end to
reverse the order of the qubits. For three qubits, that is three Hadamards, three
controlled rotations and one swap, and multiplying their matrices together gives
exactly $F_8$.

Compare the costs. The classical FFT on $N = 2^n$ numbers takes $O(N\log N) = O(2^n n)$
operations. The QFT takes $O(n^2)$ gates. That looks like an exponential speedup, and
in a narrow sense it is.

\begin{concept}{Quantum Fourier transform}
The discrete Fourier transform applied to the $2^n$ amplitudes of an $n$-qubit
state, with a circuit of $O(n^2)$ gates.

\textbf{What it fixes:} it makes periodicity hidden in an exponentially large
superposition visible to a measurement.

\textbf{What it costs:} you cannot read the spectrum. Measuring returns a single
$k$, with probability $|X_k|^2$. The QFT is not a faster FFT.
\end{concept}

\section{What you can and cannot read}

After the QFT the state is $\sum_k X_k\ket k$, and a measurement gives one $k$. If
the spectrum is flat, that $k$ says nothing. If the spectrum has a few sharp peaks,
almost every measurement lands on one of them, and a few runs are enough to infer
the spacing. The QFT is useful exactly when the spectrum is sparse: periodic
functions, or phases of eigenvalues.

This is why every claim of the form ``the QFT gives an exponential speedup for
signal processing'' should be read with suspicion. Loading a classical signal into
the amplitudes costs about as much as the classical FFT, and the output cannot be
read. The QFT pays off only when the periodic state is produced inside the
computer, as in Shor's algorithm, and only a small piece of information (a period,
a phase) needs to come out.

\section{Two uses}

\textbf{Shor.} The state $\sum_a \ket a\ket{x^a \bmod N}$ is periodic in $a$, with
period $r$. A QFT on the first register puts the probability on multiples of $q/r$,
and a measurement plus continued fractions gives $r$ (Chapter~\ref{ch:encoding}).

\textbf{Phase estimation.} A unitary $U$ with eigenvector $\ket{u}$ and eigenvalue
$e^{i2\pi\phi}$ can be used to write the phase $\phi$ into a control register, as
$\frac{1}{\sqrt N}\sum_k e^{i2\pi k\phi}\ket k$: a pure oscillation with frequency
$\phi$. The inverse QFT concentrates it on the integer closest to $N\phi$. Alexei
Kitaev introduced phase estimation in 1995, and in 1998 Cleve, Ekert, Macchiavello
and Mosca showed that Shor's algorithm is a special case of it. Phase estimation is
how a quantum computer reads eigenvalues, including the energies of a Hamiltonian,
which brings the story back to Feynman's original question.
